\section{Geometrical analysis}
Every frame, the boids have to perform a fixed-radius nearest neighbors search to determine their 
new velocities. A common solution is spatial hashing, however, we have also analysed the geometric 
properties of the boids in attempt to optimize the simulation differently. Unfortunately, these 
findings have not resulted in a practical solution, but are still worth mentioning.

The boids form a graph $G = (V, E)$, where $V$ are the boids and $E$ are the boids that are in 
viewing distance of each other. This graph changes each tick of the simulation by the positions 
in $V$ and adjusting the neighbors in $E$. Edges that stay in $E$ or are removed are easy to update; 
check the new distance and update accordingly. Finding new edges is the complex problem.

\subsection{Sweep line approach}
In a sweep line or sweep surface algorithm, a sweep line is used to sweep over data points, looping over
the data points in a lexicographical order. Data points pass the sweep line one by one and step by step
a solution is constructed for the given problem that the algorithm tackles.